\newcommand{\titulus}{\nomenFesti{Ss. Mauritii \& Sociorum Martyrum.}
\dies{Die 22. Septembris.}}
\newcommand{\oratio}{\pars{Oratio.}

\noindent Deus, qui es Sanctórum tuórum splendor mirábilis, quique hunc diem beatórum Maurítii et Sociórum eius martýrio consecrásti, da Ecclésiæ tuæ de tanta natalítii festivitáte lætári, ut apud misericórdiam tuam exémplis eórum et précibus adiuvémur.

\pars{Pro pace in universo mundo.} \scriptura{Sir. 50, 25; 2 Esdr. 4, 20; \textbf{H416}}

\vspace{-4mm}

\antiphona{II D}{temporalia/ant-dapacemdomine.gtex}

\vfill

\noindent Deus, a quo sancta desidéria, recta consília et iusta sunt ópera: da servis tuis illam, quam mundus dare non potest, pacem; ut et corda nostra mandátis tuis dédita, et hóstium subláta formídine, témpora sint tua protectióne tranquílla.

\noindent Per Dóminum nostrum Iesum Christum, Fílium tuum, qui tecum vivit et regnat in unitáte Spíritus Sancti, Deus, per ómnia sǽcula sæculórum.

\noindent \Rbardot{} Amen.}
\newcommand{\invitatorium}{\pars{Invitatorium.}

\vspace{-4mm}

\antiphona{E}{temporalia/inv-regemmartyrumsimplex.gtex}}
\newcommand{\hymnusmatutinum}{\pars{Hymnus}

{
\grechangedim{interwordspacetext}{0.16 cm plus 0.15 cm minus 0.05cm}{scalable}%
\cuminitiali{VIII}{temporalia/hym-MartyrumRector.gtex}
\grechangedim{interwordspacetext}{0.22 cm plus 0.15 cm minus 0.05cm}{scalable}%
}
}
\newcommand{\lectioi}{\pars{Lectio I.} \scriptura{Idt. 6, 1-3.5. 10-21}

\noindent De libro Iudith.

\noindent Et, cum cessásset tumúltus virórum, qui erant in circúitu concílii, dixit Holoférnes princeps milítiæ virtútis Assýriæ ad Achior coram omni pópulo alienigenárum et ad omnes fílios Moab: “ Et quis es tu, Achior, et mercennárii Ephraim, quóniam prophetásti nobis sicut hódie et dixísti quóniam genus filiórum Israel non expugnátur, quóniam Deus eórum próteget eos? Et quis est Deus, nisi Nabuchodónosor rex omnis terræ? Hic mittet potestátem suam et dispérdet eos a fácie terræ, et non liberábit eos Deus eórum. Sed nos servi illíus percutiémus eos sicut hóminem unum, nec sustinébunt vim equórum nostrórum. Tu autem, Achior, mercennárii Ammon, qui locútus es verba hæc in die iniquitátis tuæ, non vidébis fáciem meam iam ex hac die, quoadúsque ulcíscar in genus illórum, qui ascendérunt ex Ægýpto.

\noindent Et præcépit Holoférnes servis suis, qui erant adstántes in tabernáculo eius, ut comprehénderent Achior et dúcerent eum in Betúliam et tráderent in manus filiórum Israel. Et comprehendérunt eum servi eius et duxérunt illum extra castra in campum et promovérunt ex médio campo in montánam et advenérunt ad fontes, qui erant sub Betúlia. Et, ut vidérunt eos viri civitátis super vérticem montis, accepérunt arma sua et exiérunt extra civitátem supra cacúmen montis, et omnis vir fundibalárius occupávit ascénsus eórum et mittébat lápides super eos. Et accedéntes sub monte alligavérunt Achior et reliquérunt proiéctum sub radíce montis et revérsi sunt ad dóminum suum.}
\newcommand{\responsoriumi}{\pars{Responsorium 1.} \scriptura{\Rbardot{} Idt. 9, 10 \Vbardot{} ibid. 9, 11; \textbf{H410}}

\vspace{-5mm}

\responsorium{II}{temporalia/resp-dominedeusquiconteris-CROCHU.gtex}{}}
\newcommand{\lectioii}{\pars{Lectio II.}

\noindent Ex libro Sancti Euchérii Epíscopi de Passióne Agaunénsium Mártyrum.

\noindent Sub Maximiáno, qui Románæ rei públicæ cum Diocletiáno colléga impérium ténuit, per divérsas fere províncias laniáti aut interfécti sunt mártyrum pópuli. Erat eódem témpore in exércitu légio mílitum, qui Thebǽi appellabántur. Hi in auxílium Maximiáno ab Oriéntis pártibus accíti vénerant, viri in rebus béllicis strénui et virtúte nóbiles, sed nobilióres fide; erga imperatórem fortitúdine, erga Christum devotióne certábant. Evangélici præcépti étiam sub armis non immémores reddébant quæ Dei erant Deo, et quæ Cǽsaris Cǽsari restituébant.

\noindent Itaque cum ad pertrahéndam Christianórum multitúdinem destinaréntur, soli crudelitátis ministérium detrectáre ausi sunt atque huiúsmodi præcéptis se obtemperatúros negant. Maximiánus, qui se circa Octodúrum itínere fessus tenébat, ubi cum sibi per núntios delátum esset Legiónem hanc in Agaunénsibus angústiis substitísse, in furórem instínctu indignatiónis exársit. Cógnito ígitur Thebæórum respónso, Maximiánus, præcípiti ira férvidus, décimum quemque gládio feríri iubet, quo facílius céteri, régiis præcéptis térriti, metu céderent; redintegratísque mandátis edícit, ut réliqui in persecutiónem Christianórum cogántur.}
\newcommand{\responsoriumii}{\pars{Responsorium 2.} \scriptura{\Vbardot{} Mt. 25, 34; \textbf{H331}}

\vspace{-5mm}

\responsorium{VIII}{temporalia/resp-sanctimei-CROCHU-sinedox.gtex}{}

\rubrica{vel ad libitum:}

\vspace{3mm}

\pars{Responsorium 2.} \scriptura{\Vbardot{} Sap. 5, 6; \textbf{H369}}

\vspace{-5mm}

\responsorium{VII}{temporalia/resp-fulgebuntiusti-CROCHU.gtex}{}}
\newcommand{\lectioiii}{\pars{Lectio III.}

\noindent Quibus iussis dénuo in castra perlátis, segregátus atque percússus est, qui décimus sorte obvénerat, réliqua vero se mílitum multitúdo mútuo sermóne instigábat, ut in tam præcláro ópere persísteret.

\noindent Thebǽi Maximiáno mandáta hæc mittunt: Mílites sumus, imperátor, tui, sed tamen servi, quod líbere confitémur, Dei; tibi milítiam debémus, illi innocéntiam. Offérimus nostras in quémlibet hostem manus, quas sánguine innocéntium cruentáre nefas dúcimus. Iurávimus primum in sacraménta divína, iurávimus deínde in sacraménta régia; nihil nobis de secúndis credas necésse est, si prima perrúmpimus. Christiános ad pœnam per nos requíri iubes; iam tibi ex hoc álii requiréndi non sunt: Christiános nos fatémur, pérsequi Christiános non póssumus.

\noindent Maximiánus, obstinátos in fide Christi cernens ánimos virórum, desperánsque gloriósam eórum constántiam posse revocári, una senténtia intérfici omnes decrévit et rem cónfici circumfúsis mílitum agmínibus iubet.}
\newcommand{\responsoriumiii}{\pars{Responsorium 3.} \scriptura{\Vbardot{} Ap. 7, 14; \textbf{H362}}

\vspace{-5mm}

\responsorium{IV}{temporalia/resp-istisunttriumphatores.gtex}{}}
\newcommand{\laudes}{\pars{Psalmus 1.} \scriptura{\textbf{H311}}

\antiphona{VIII G}{temporalia/ant-sanctusmauritius.gtex}

%\vspace{-2mm}

\scriptura{Psalmus 62}

%\vspace{-2mm}

\initiumpsalmi{temporalia/ps62-initium-viii-G-auto.gtex}

%\vspace{-1.5mm}

\input{temporalia/ps62-viii-G.tex} \Abardot{}

\vfill
\pagebreak

\pars{Psalmus 2.} \scriptura{\textbf{H311}}

\antiphona{III b}{temporalia/ant-sanctalegioagaunensiummartyrum.gtex}

%\vspace{-2mm}

\scriptura{Canticum trium puerorum, Dan. 3, 57-88 et 56}

\initiumpsalmi{temporalia/dan3-initium-iii-b-auto.gtex}

\input{temporalia/dan3-iii-b-sinedox.tex}

\rubrica{Hic non dicitur Gloria Patri, neque Amen.}

\vfill

\antiphona{}{temporalia/ant-sanctalegioagaunensiummartyrum.gtex}

\vfill
\pagebreak

\pars{Psalmus 3.} \scriptura{\textbf{H311}}

\vspace{-4mm}

\antiphona{III a\textsuperscript{2}}{temporalia/ant-pretiosasunt.gtex}

%\vspace{-2mm}

\scriptura{Psalmus 149}

%\vspace{-2mm}

\initiumpsalmi{temporalia/ps149-initium-iii-a2-auto.gtex}

\input{temporalia/ps149-iii-a2.tex} \Abardot{}

\vfill
\pagebreak}
\newcommand{\lectiobrevis}{\pars{Lectio Brevis.} \scriptura{Ap. 7, 14-15}

\noindent Hi sunt qui venérunt de tribulatióne magna, et lavérunt stolas suas, et dealbavérunt eas in sánguine Agni. Ideo sunt ante thronum Dei, et sérviunt ei die ac nocte in templo eius.}
\newcommand{\responsoriumbreve}{\pars{Responsorium breve.} \scriptura{Sap. 5, 16}

\antiphona{VI}{temporalia/resp-iustiautem.gtex}}
\newcommand{\preces}{\noindent Fratres, Salvatórem nostrum, testem fidélem, per mártyres interféctos propter verbum Dei, celebrémus,~\grestar{} clamántes:

\Rbardot{} Redemísti nos Deo in sánguine tuo.

\noindent Per mártyres tuos, qui líbere mortem in testimónium fídei sunt ampléxi,~\grestar{} da nobis, Dómine, veram spíritus libertátem.

\Rbardot{} Redemísti nos Deo in sánguine tuo.

\noindent Per mártyres tuos, qui fidem usque ad sánguinem sunt conféssi,~\grestar{} da nobis, Dómine, puritátem fideíque constántiam.

\Rbardot{} Redemísti nos Deo in sánguine tuo.

\noindent Per mártyres tuos, qui, sustinéntes crucem, tua vestígia sunt secúti,~\grestar{} da nobis, Dómine, ærúmnas vitæ fórtiter sustinére.

\Rbardot{} Redemísti nos Deo in sánguine tuo.

\noindent Per mártyres tuos, qui stolas suas lavérunt in sánguine Agni,~\grestar{} da nobis, Dómine, omnes insídias carnis mundíque devíncere.

\Rbardot{} Redemísti nos Deo in sánguine tuo.}
\newcommand{\benedictus}{\pars{Canticum Zachariæ.} \scriptura{\textbf{H311}}

\vspace{-4mm}

\antiphona{I D*}{temporalia/ant-flagrabatinbeatissimis.gtex}

%\vspace{-2mm}

\scriptura{Lc. 1, 68-79}

%\vspace{-2mm}

\cantusSineNeumas
\initiumpsalmi{temporalia/benedictus-initium-i-D_-auto.gtex}

%\vspace{-1.5mm}

\input{temporalia/benedictus-i-D_.tex} \Abardot{}

\antiphona{}{temporalia/ant-flagrabatinbeatissimis.gtex}}
\newcommand{\benedicamuslaudes}{\cuminitiali{}{temporalia/benedicamus-memoria-laudes.gtex}}
\include{hebdomadaxxv}
\include{feriaiii}
